\documentclass[11pt]{article}
\usepackage{amsmath, amssymb}
\usepackage{geometry}
\geometry{margin=1in}

\title{Momentum-Based Multi-Robot Control:\\A First-Order Actuator Perspective}
\author{}
\date{}

\begin{document}
\maketitle

\section{System Model}

\subsection{Robot Hardware}
We consider $n \geq 3$ mobile robots operating on a 2D plane. Each robot accepts velocity commands $\mathbf{v}_i = [v_{x,i}, v_{y,i}, \omega_i]^T$ and possesses an internal velocity controller that tracks commands with a first-order response. 

When a step velocity command of 0.3~m/s is applied from rest, the robot's actual velocity follows:
\begin{equation}
v(t) = v_{\text{cmd}}(1 - e^{-t/\tau})
\end{equation}
where $\tau$ is the time constant characterizing the actuator and its built-in controller.

\subsection{Rigid Formation Assumption}
We assume the robots maintain a rigid formation through a high-bandwidth formation controller. This allows us to model the entire cluster as a single rigid body with translational velocity $\dot{\mathbf{x}}$ at the geometric center and rotational velocity $\omega$ about that center.

\subsection{Control Architecture}
Our system has two layers:
\begin{enumerate}
    \item \textbf{Layer 1 -- Robot Inner Loop:} Each robot's built-in controller converts commanded velocity $\mathbf{v}_{\text{cmd}}$ into actual velocity $\mathbf{v}$ through a first-order low-pass response with time constant $\tau$.
    
    \item \textbf{Layer 2 -- Formation Controller:} Our controller computes desired cluster velocity from position error:
    \begin{equation}
    \mathbf{v}_{\text{des}}(t) = k_v(\mathbf{x}_{\text{des}}(t) - \mathbf{x}(t))
    \end{equation}
    where $k_v$ [Hz] is a proportional gain mapping position error to target velocity.
\end{enumerate}

The momentum behavior emerges directly from the first-order nature of Layer 1.

\section{Continuous-Time Dynamics}

The robot's actuator behaves as a first-order system:
\begin{equation}
\dot{\mathbf{v}}(t) = -\frac{1}{\tau}\mathbf{v}(t) + \frac{1}{\tau}\mathbf{v}_{\text{des}}(t)
\label{eq:continuous}
\end{equation}

This is a low-pass filter: for constant commands, it produces exponential convergence to steady-state; for time-varying commands, it attenuates high-frequency changes. The time constant $\tau$ is determined purely by the actuator hardware—no virtual mass or damping is assumed.

\section{Discrete-Time Implementation}

\subsection{Discretization}
For a digital controller with period $\Delta t$, we discretize Eq.~\eqref{eq:continuous}. Integrating the continuous dynamics over one time step yields:
\begin{equation}
\mathbf{v}[k+1] = \alpha \mathbf{v}[k] + (1-\alpha)\mathbf{v}_{\text{des}}[k]
\label{eq:momentum}
\end{equation}
where the momentum coefficient is:
\begin{equation}
\alpha = e^{-\Delta t/\tau}
\end{equation}

This is the fundamental momentum equation. The parameter $\alpha \in (0,1)$ characterizes the actuator's physical response:
\begin{itemize}
    \item $\alpha \to 1$: Sluggish actuator, strong momentum retention (large $\tau$)
    \item $\alpha \to 0$: Fast actuator, minimal momentum (small $\tau$)
\end{itemize}

\subsection{Complete Control Law}
Substituting the formation controller's desired velocity into Eq.~\eqref{eq:momentum}, the final commanded velocity becomes:

\begin{equation}
\boxed{\mathbf{v}_{\text{cmd}}[k+1] = \alpha \mathbf{v}_{\text{cmd}}[k] + (1-\alpha)k_v(\mathbf{x}_{\text{des}}[k] - \mathbf{x}[k])}
\label{eq:final}
\end{equation}

The same form applies to the rotational channel $\omega$. This controller has two intuitive parameters:
\begin{itemize}
    \item $k_v$: Controls how aggressively position error drives desired velocity
    \item $\alpha$: Fixed by actuator dynamics (or increased for additional smoothing)
\end{itemize}

\section{Parameter Identification}

\subsection{Experimental Procedure}
To determine $\alpha$ from robot hardware:
\begin{enumerate}
    \item Start a robot at rest
    \item Apply a step velocity command (e.g., 0 $\to$ 0.3 m/s)
    \item Record the velocity response over time
    \item Fit to $v(t) = v_{\text{cmd}}(1 - e^{-t/\tau})$ to extract $\tau$
    \item Calculate $\alpha = e^{-\Delta t/\tau}$ for your control rate
\end{enumerate}

\subsection{Example}
For a 10~Hz control loop ($\Delta t = 0.1$~s) with $\tau = 0.3$~s:
\begin{equation}
\alpha = e^{-0.1/0.3} = e^{-0.333} \approx 0.717
\end{equation}

Under the rigid formation assumption, this single measurement characterizes the entire cluster's momentum behavior.

\end{document}